\scanpage{12}
\mspage{10}

Toutes ces d\unsure{es}criptions sont compatibles avec les opérations tensorielles, et la formation d'images inverses pour k variable.

On peut se demander, avec la description précédente des bi-isocristaux, quels sont ceux qui sont dans l'image essentielle de $\unsure{\operatorname{Isbicr}|\operatorname{Hsbicr}}(S,i)\longrightarrow\operatorname{Isbicr}(S)$. On trouve que ce sont les couples $(E,F)$, E un vectoriel sur K et F un automorphisme semi-li\unsure{né}aire, tels que, si on pose $V=p^iF^{-1}$, pour tout $x\in E$, l'ensemble de ses transformés par les différents monômes non commutatifs en F,V soit une partie \uline{bornée} de E: c'est une pure tautologie. Pour qu'il existe un i ayant cette propriété, il faut et il suffit que pour tout $x\in E$, l'ensemble des $F^nx$ soit borné. Bien entendu, c'est là une propriété qui est stable par changement de F en pF (correspondant à la tensorisation par le bi-isocristal de Tate de poids 2), mais non par le changement de F en $p^{-1}F$ (correspondant à la tensorisa\unsure{tion|gion} par l'inverse $T^{-1}$ du bi-isocristal précédent). Comme on tient beaucoup à inverser Tate, on voit qu'il n'est pas naturel, en effet, de poser une condition de croissance sur l'ensemble des \deleted{\unsure{iétér|itér}} itérés de F. De façon précise, appelant \uline{effectifs} les bi-\unsure{is}ocristaux qui proviennent d'un objet d'un $\operatorname{Isbicr}(S,i)$, on voit que tout bi-isocristal est de la forme $T^{-i}\otimes\underline N$, avec $\underline N$ effectif: on n'a pas ajouté plus de nouveaux objets qu'il n'était absolument nécessaire !

1.13. \uline{Exemple} 2: S \uline{lisse sur un corps parfait} k. Déterminons d'abord dans ce cas les cristaux sur S, sans plus. \unsure{Tout|Wout} d'abord, sans condition de lissité, on voit à l'aide de la propriété caractéristique de W que la catégorie C introduite dans 1.0. ne change pas, essentiellement, quand on remplace $R=\operatorname{Spec}(\underline Z)$ par $R=\operatorname{Spec}(W)$. D'autre part, il résulte aussitôt des définitions que la donnée d'un \deleted{\unsure{R-}}cristal sur S, relativement à R, revient à la donnée d'un système cohérent de cristaux sur S, relatifs aux $W_n=W/p^{n+1}W$. (NB On n'a pas formulé avec la généralité qui convenait l'exemple 3 de 1.7. ! ) Théoriquement, on est donc ramené à déterminer, pour chaque n, la catégorie des cristaux sur S relatifs à $W_n$, et les foncteurs restriction entre ces catégories. D'ailleurs, il est inoffensif de se localiser sur S, \deleted{\unsure{Ceci|aussi}} par exemple de supposer au besoin S affine. Utilisant maintenant la lissité de S, on peut donc supposer que pour tout n, S se remonte en
\margin{11}
